% Module 0: Introduction and Motivation
% Estimated slides: 20 | Duration: ~1 hour
% Audience: Engineers/programmers
% Key topics: Why QC, classical limits, applications, history, industry landscape, course map
% Labs/Exercises: Industry use-case quiz

%%%%%%%%%%%%%%%%%%%%%%%%%%%%%%%%%%%%%%%%%%%%%%%%%%%%%%%%%%
\begin{frame}[fragile]\frametitle{}
\begin{center}
{\Large Quantum Computing for Non-Physicists}\\[0.5em]
{\large Module 0: Introduction and Motivation}
\end{center}
\end{frame}

%%%%%%%%%%%%%%%%%%%%%%%%%%%%%%%%%%%%%%%%%%%%%%%%%%%%%%%%%%%
\begin{frame}[fragile]\frametitle{What This Course Is (and Is Not)}
\begin{itemize}
\item This is NOT a physics course --- no wave functions, no Schrodinger equation
\item This IS a programmer's guide to quantum computing
\item You will understand what quantum computers do and why it matters
\item You will write and run real quantum circuits using Qiskit
\item Assumed background: school-level physics, programming in Python
\end{itemize}
% Suggested diagram: Course roadmap visual with modules as nodes
\end{frame}

%%%%%%%%%%%%%%%%%%%%%%%%%%%%%%%%%%%%%%%%%%%%%%%%%%%%%%%%%%%
\begin{frame}[fragile]\frametitle{Quantum Computing Fundamentals}
\begin{itemize}
\item Introduction to quantum computing concepts
\item From classical bits to quantum qubits
\item Quantum gates and circuits
\item Applications and future prospects
\item Based on educational resources and industry insights
\end{itemize}
\end{frame}

%%%%%%%%%%%%%%%%%%%%%%%%%%%%%%%%%%%%%%%%%%%%%%%%%%%%%%%%%%%
\begin{frame}[fragile]\frametitle{Why Should Programmers Care?}
\begin{itemize}
\item Classical computers are hitting physical limits (transistors approaching atom size)
\item Quantum computers solve certain problems exponentially faster
\item Drug discovery: simulating molecules is intractable classically for large molecules
\item Cryptography: current encryption will be broken by sufficiently large quantum computers
\item Optimization: logistics, finance, machine learning training
\item Industry is investing billions --- IBM, Google, Microsoft, Amazon, IonQ
\end{itemize}
% Suggested diagram: Bar chart of quantum investment by year 2015-2025
\end{frame}

%%%%%%%%%%%%%%%%%%%%%%%%%%%%%%%%%%%%%%%%%%%%%%%%%%%%%%%%%%%
\begin{frame}[fragile]\frametitle{Introduction: Why Quantum Computing Matters}
\begin{itemize}
\item Many unsolved problems limited by computation, not data
\item Examples: drug discovery, supply chains, AI training, optimization
\item Classical computers hit limits for certain problem categories
\item Quantum computing offers fundamentally different computational model
\item New ways of searching and optimizing
\item Ability to solve problems unreachable for classical systems
\end{itemize}
\end{frame}

%%%%%%%%%%%%%%%%%%%%%%%%%%%%%%%%%%%%%%%%%%%%%%%%%%%%%%%%%%%
\begin{frame}[fragile]\frametitle{The Classical Computing Wall}
\begin{itemize}
\item Moore's Law: transistor count doubles every ~2 years (since 1965)
\item Transistors are now just a few atoms wide --- quantum effects become a problem
\item Clock speed has stalled around 3-4 GHz since 2004
\item More cores help, but some problems cannot be parallelized classically
\item Energy cost of computation is rising unsustainably
\end{itemize}
% Suggested diagram: Moore's Law graph showing transistor count vs year, plateau in clock speed
\end{frame}

%%%%%%%%%%%%%%%%%%%%%%%%%%%%%%%%%%%%%%%%%%%%%%%%%%%%%%%%%%%
\begin{frame}[fragile]\frametitle{NP-Hard Problems}
\begin{itemize}
\item Problems where complexity grows exponentially
\item Each additional variable dramatically increases computation
\item Example: 1 variable = 2 possibilities, 10 variables = 1024 possibilities
\item 100 variables = astronomically large search space
\item Creates combinatorial explosion
\item Resource requirements grow as $2^n$
\item Classical computers eventually run out of memory and time
\end{itemize}
\end{frame}

%%%%%%%%%%%%%%%%%%%%%%%%%%%%%%%%%%%%%%%%%%%%%%%%%%%%%%%%%%%
\begin{frame}[fragile]\frametitle{Why Quantum Helps}
\begin{itemize}
\item Exploits superposition and quantum state combinations
\item Parallel probability exploration
\item Some exponential problems behave more like linear
\item Instead of checking one state at a time, quantum systems represent multiple states
\item Massive search advantages
\item Optimization advantages
\item Potential exponential speedups for certain problems
\end{itemize}
\end{frame}

%%%%%%%%%%%%%%%%%%%%%%%%%%%%%%%%%%%%%%%%%%%%%%%%%%%%%%%%%%%
\begin{frame}[fragile]\frametitle{What Makes Quantum Different?}
\begin{itemize}
\item Quantum computers exploit quantum mechanical properties of matter
\item A quantum bit (qubit) can be 0, 1, or both at the same time (superposition)
\item Multiple qubits can be linked in ways classical bits cannot (entanglement)
\item These properties let quantum computers explore many solutions simultaneously
\item Not faster at everything --- quantum advantage exists for specific problem types
\end{itemize}
% Suggested diagram: Classical bit (0 or 1) vs qubit (sphere with arrow) side by side
\end{frame}

%%%%%%%%%%%%%%%%%%%%%%%%%%%%%%%%%%%%%%%%%%%%%%%%%%%%%%%%%%%
\begin{frame}[fragile]\frametitle{Classical vs Quantum Computing}
\begin{itemize}
\item Classical computers excel at: spreadsheets, web browsing, standard software
\item Handle linear or polynomial complexity problems well
\item Linear scaling: 10x workload requires roughly 10x resources
\item Quantum computers are NOT replacements for classical
\item They are specialized accelerators for specific problem classes
\item Analogy: GPUs did not replace CPUs, they complement them
\item Quantum computers will coexist with classical systems
\end{itemize}
\end{frame}

%%%%%%%%%%%%%%%%%%%%%%%%%%%%%%%%%%%%%%%%%%%%%%%%%%%%%%%%%%%
\begin{frame}[fragile]\frametitle{Key Applications of Quantum Computing}
\begin{itemize}
\item \textbf{Drug Discovery}: Simulate molecular interactions at quantum level
\item \textbf{Cryptography}: Break RSA encryption; quantum-safe replacements
\item \textbf{Optimization}: Portfolio optimization, logistics routing, scheduling
\item \textbf{Machine Learning}: Faster training, quantum kernels, quantum neural networks
\item \textbf{Material Science}: Design new superconductors, batteries, catalysts
\item \textbf{Weather/Climate}: High-fidelity climate simulations
\end{itemize}
% Suggested diagram: 6-panel icons for each application domain
\end{frame}

%%%%%%%%%%%%%%%%%%%%%%%%%%%%%%%%%%%%%%%%%%%%%%%%%%%%%%%%%%%
\begin{frame}[fragile]\frametitle{Shor's Algorithm \& Encryption}
\begin{itemize}
\item One of the most famous quantum algorithms
\item Modern encryption depends on prime factorization difficulty
\item RSA encryption: $N = p \times q$ (two large primes)
\item Easy to multiply primes, extremely hard to factor
\item Classical computers struggle with factorization
\item Shor's algorithm: quantum computers can factor quickly
\item Implication: current encryption systems could break
\item Post-quantum cryptography becoming important
\end{itemize}
\end{frame}

%%%%%%%%%%%%%%%%%%%%%%%%%%%%%%%%%%%%%%%%%%%%%%%%%%%%%%%%%%%
\begin{frame}[fragile]\frametitle{A Brief History of Quantum Computing}
\begin{itemize}
\item 1980: Feynman proposes using quantum systems to simulate quantum physics
\item 1994: Shor's algorithm -- quantum factoring threatens RSA encryption
\item 1996: Grover's algorithm -- quadratic speedup for search
\item 2001: IBM demonstrates 7-qubit Shor's algorithm
\item 2016: IBM puts a 5-qubit quantum computer on the cloud (IBM Quantum Experience)
\item 2019: Google claims ``quantum supremacy'' with 53-qubit Sycamore processor
\item 2023+: IBM 1000+ qubit processors; IonQ, Quantinuum, PsiQuantum scale up
\end{itemize}
% Suggested diagram: Timeline with key milestones
\end{frame}

%%%%%%%%%%%%%%%%%%%%%%%%%%%%%%%%%%%%%%%%%%%%%%%%%%%%%%%%%%%
\begin{frame}[fragile]\frametitle{The Quantum Computing Landscape Today}
\begin{itemize}
\item \textbf{Superconducting}: IBM (Eagle, Osprey, Heron), Google (Sycamore), Rigetti
\item \textbf{Trapped Ion}: IonQ, Honeywell/Quantinuum
\item \textbf{Photonic}: PsiQuantum, Xanadu
\item \textbf{Neutral Atom}: Pasqal, QuEra, Atom Computing
\item \textbf{Silicon Spin}: Intel, Silicon Quantum Computing
\item \textbf{Cloud access}: IBM Quantum, AWS Braket, Azure Quantum, Google Cloud Quantum AI
\end{itemize}
% Suggested diagram: Technology landscape bubble chart (company vs technology type)
\end{frame}

%%%%%%%%%%%%%%%%%%%%%%%%%%%%%%%%%%%%%%%%%%%%%%%%%%%%%%%%%%%
\begin{frame}[fragile]\frametitle{Quantum Industry Landscape}
\begin{itemize}
\item Major players: IBM, Google, Microsoft, and others
\item Different architectural approaches being explored
\item Superconducting qubits (IBM, Google)
\item Photonic systems (PsiQuantum, Xanadu)
\item Trapped ions (IonQ, Honeywell/Quantinuum)
\item Neutral atoms, topological approaches
\item No single dominant architecture yet
\end{itemize}
\end{frame}

%%%%%%%%%%%%%%%%%%%%%%%%%%%%%%%%%%%%%%%%%%%%%%%%%%%%%%%%%%%
\begin{frame}[fragile]\frametitle{NISQ Era: Where We Are Now}
\begin{itemize}
\item NISQ = Noisy Intermediate-Scale Quantum
\item Today's machines: 50-1000+ physical qubits, but noisy (error rates ~0.1-1\%)
\item Not yet fault-tolerant --- computations are limited in depth
\item Still useful for certain problems: chemistry simulation, optimization, ML
\item Fault-tolerant quantum computing (FTQC) is the long-term goal
\item FTQC may need millions of physical qubits for useful logical qubits
\end{itemize}
% Suggested diagram: Road to FTQC showing NISQ -> error-corrected -> fault-tolerant stages
\end{frame}

%%%%%%%%%%%%%%%%%%%%%%%%%%%%%%%%%%%%%%%%%%%%%%%%%%%%%%%%%%%
\begin{frame}[fragile]\frametitle{The Quantum Software Stack}
\begin{itemize}
\item \textbf{Application layer}: Algorithms (Grover, Shor, VQE, QAOA)
\item \textbf{Framework layer}: Qiskit (IBM), Cirq (Google), PennyLane, Braket (Amazon)
\item \textbf{Compiler/transpiler}: Converts circuits to native gate sets
\item \textbf{Hardware layer}: Actual qubits (superconducting, trapped ion, etc.)
\item As a programmer, you mostly work at the framework and algorithm layer
\end{itemize}
% Suggested diagram: Stack diagram with layers
\end{frame}

%%%%%%%%%%%%%%%%%%%%%%%%%%%%%%%%%%%%%%%%%%%%%%%%%%%%%%%%%%%
\begin{frame}[fragile]\frametitle{How to Think About Quantum Computing}
\begin{itemize}
\item Think of qubits as \textit{probabilistic} bits, not just 0 or 1
\item Quantum algorithms set up probability amplitudes so right answers are likely
\item Measurement ``collapses'' the quantum state --- run circuits many times
\item Intuition, not equations, is your primary tool as a programmer
\item Classical analogies help initially but will break down --- embrace the weirdness
\end{itemize}
\end{frame}

%%%%%%%%%%%%%%%%%%%%%%%%%%%%%%%%%%%%%%%%%%%%%%%%%%%%%%%%%%%
\begin{frame}[fragile]\frametitle{Course Roadmap}
\begin{itemize}
\item \textbf{Workshop A}: Bits to Qubits, Rules of Quantum World (Foundations)
\item \textbf{Workshop B}: Qubit Gates, Circuits, Qiskit Programming
\item \textbf{Workshop C}: Quantum Algorithms (Grover, QFT, Shor's, VQE, QAOA)
\item \textbf{Workshop D}: Hardware, Communication, Information, Simulations, Applications
\item Throughout: Hands-on Qiskit labs after every major topic
\item Capstone: Mini-project on a quantum computing problem of your choice
\end{itemize}
% Suggested diagram: Course map with 4 workshops and milestone labs
\end{frame}

%%%%%%%%%%%%%%%%%%%%%%%%%%%%%%%%%%%%%%%%%%%%%%%%%%%%%%%%%%%
\begin{frame}[fragile]\frametitle{Tools and Setup}
\begin{itemize}
\item \textbf{Python 3.8+} and \textbf{Jupyter Notebook} (Anaconda recommended)
\item \textbf{Qiskit}: \texttt{pip install qiskit qiskit-aer qiskit-ibm-runtime}
\item \textbf{IBM Quantum account}: Free at \texttt{quantum.ibm.com} --- needed for real hardware
\item \textbf{Optional}: PennyLane (\texttt{pip install pennylane}) for QML demos
\item Verify: \texttt{import qiskit; print(qiskit.version.get\_version\_info())}
\end{itemize}
% Suggested diagram: Installation checklist screenshot
\end{frame}

%%%%%%%%%%%%%%%%%%%%%%%%%%%%%%%%%%%%%%%%%%%%%%%%%%%%%%%%%%%
\begin{frame}[fragile]\frametitle{Quantum Computing Is NOT Faster Excel}
\begin{itemize}
\item NOT: faster Excel, browsing, or Zoom calls
\item IS: specialized compute for specific problems
\item Optimization engine for complex decision problems
\item Search accelerator for large solution spaces
\item Scientific simulation system for quantum phenomena
\item Complements classical computing infrastructure
\item Think of it as a specialized coprocessor
\end{itemize}
\end{frame}

%%%%%%%%%%%%%%%%%%%%%%%%%%%%%%%%%%%%%%%%%%%%%%%%%%%%%%%%%%%
\begin{frame}[fragile]\frametitle{What Quantum Computers Cannot Do}
\begin{itemize}
\item They are NOT general-purpose speedup machines for all problems
\item They will NOT replace classical computers --- they work alongside them
\item They are NOT faster for most everyday tasks (web, video, databases)
\item Quantum advantage requires the problem to have exploitable quantum structure
\item Understanding limitations is as important as understanding capabilities
\end{itemize}
\end{frame}

%%%%%%%%%%%%%%%%%%%%%%%%%%%%%%%%%%%%%%%%%%%%%%%%%%%%%%%%%%%
\begin{frame}[fragile]\frametitle{Exercise: Quantum Use Case Quiz}
\begin{itemize}
\item Which of these problems might benefit from quantum computing? (discuss in pairs)
\begin{enumerate}
\item Finding a name in a sorted phone book
\item Factoring a 2048-bit number used in RSA encryption
\item Simulating the folding of a protein with 50 amino acids
\item Sorting 1 million integers
\item Finding the optimal route for 100 delivery vehicles
\item Training a deep neural network on ImageNet
\end{enumerate}
\item Answers and reasoning: next slide
\end{itemize}
\end{frame}

%%%%%%%%%%%%%%%%%%%%%%%%%%%%%%%%%%%%%%%%%%%%%%%%%%%%%%%%%%%
\begin{frame}[fragile]\frametitle{Exercise: Quantum Use Case Quiz (Answers)}
\begin{itemize}
\item (1) Sorted search: classical binary search is already O(log n) --- no quantum advantage
\item (2) RSA factoring: YES --- Shor's algorithm gives exponential speedup
\item (3) Protein folding: YES --- quantum simulation is exponentially more efficient
\item (4) Sorting: No clear quantum advantage (Grover gives at best $O(\sqrt{n})$ comparisons)
\item (5) Vehicle routing: YES --- QAOA can tackle combinatorial optimization
\item (6) Neural network training: Unclear; quantum ML is active research, not proven yet
\end{itemize}
\end{frame}

%%%%%%%%%%%%%%%%%%%%%%%%%%%%%%%%%%%%%%%%%%%%%%%%%%%%%%%%%%%
\begin{frame}[fragile]\frametitle{Important Terms}
\begin{itemize}
\item Qubit: quantum version of a bit
\item Superposition: ability to exist in multiple states simultaneously
\item NP-hard: problems with exponential complexity
\item Optimization: finding best solution among many possibilities
\item QUBO: Quadratic Unconstrained Binary Optimization format
\item Ising Model: physics-based optimization framework
\item Shor's Algorithm: quantum factorization algorithm
\item Photonic Quantum Computing: quantum systems using photons/light
\end{itemize}
\end{frame}

%%%%%%%%%%%%%%%%%%%%%%%%%%%%%%%%%%%%%%%%%%%%%%%%%%%%%%%%%%%
\begin{frame}[fragile]\frametitle{Reading List and Resources}
\begin{itemize}
\item \textbf{Books}: ``Quantum Computing: An Applied Approach'' by Jack Hidary; ``Programming Quantum Computers'' by Johnston et al.
\item \textbf{Online}: IBM Quantum Learning (learning.quantum.ibm.com); Qiskit Textbook (qiskit.org/learn)
\item \textbf{Courses}: MIT 8.370x; Berkeley CS 294; Coursera ``Quantum Computing from Basics to Cutting Edge''
\item \textbf{YouTube}: Qiskit channel, Quantum Sense, 3Blue1Brown (linear algebra)
\item \textbf{Community}: Qiskit Slack, Quantum Computing Stack Exchange
\end{itemize}
\end{frame}
